\makeatletter
\def\input@path{{./}{../}{../../}{../../../}{../../../../}{preamble/}{../preamble/}{../../preamble/}{../../../preamble/}}
\makeatother

\input{preamble_exams.tex}
\usepackage{needspace}
\usepackage{titling}
\setlength{\droptitle}{-2em}
\pretitle{\par\begingroup\centering\LARGE}
\posttitle{\par\endgroup}
\preauthor{}\postauthor{}
\predate{}\postdate{}
\maketitlehookd{\vspace{-1em}}
\linespread{1.12}
\setlength{\parskip}{5pt}
\newgeometry{margin=0.65in,top=0.55in,bottom=0.65in}

\title{Grade 11 Term 2 Home Practice -- C9\\\large Checking confidence-interval inferences\\[0.25em]
	\normalsize Global Economics -- 2026--2027}
\author{}
\date{}

\begin{document}
	\maketitle

	\setcounter{tocdepth}{1}
	\setcounter{secnumdepth}{2}
	\phantomsection
	\label{toc}
	\tableofcontents

	\section*{Evaluated criteria}
	\begin{enumerate}[label=\textbf{C\arabic*},start=9]
		\item Assesses the validity and limitations of confidence-interval inferences by evaluating the sampling design, potential bias, and independence, including the 10\% condition for sampling without replacement and procedure-specific requirements concerning sample size, distributional shape, and outliers.
	\end{enumerate}

	\textbf{How to practise.} Try each task before opening its blue solution link. The arrows return to the contents, problem list or current problem.

	\section{Worked Solution Examples}
	\label{sec:worked-solution-examples}

	\SectionProblemsTableOfContents{
		\SectionProblemLink{prob:c9-1}{A bank's card-spending report (C9)} \\
		\SectionProblemLink{prob:c9-2}{Rare defaults in a loan audit (C9)} \\
		\SectionProblemLink{prob:c9-3}{An unusual charge in a small import audit (C9)}
	}

	% ----------------------------------------------------------------
	\subsection{A bank's card-spending report (C9)}
	\label{prob:c9-1}

	A bank wants to infer mean monthly card spending for all 600 new digital-account customers. Its app shows an optional spending survey to customers who open a budgeting feature. The first 120 volunteers respond. The bank proposes a one-sample $t$ interval. The responses are strongly right-skewed and one very high spending value has not been verified.

	\QuestionsTasks[sec:worked-solution-examples]{
		{9}/{sol:c9-1-a}/
		{Identify the target population and two ways the respondent group could differ from it. Explain whether using more volunteers solves either issue.},
		{9}/{sol:c9-1-b}/
		{Check what the 10\% condition would say if 120 distinct customers had been sampled at random without replacement. Assess whether $n=120$ settles the concern about shape or outliers. Propose a defensible redesign.}
	}

	\clearpage

	\Needspace{5\baselineskip}
	\phantomsection
	\label{sol:c9-1-a}
	\subsubsection*{C9 -- Audit selection and possible bias}

	The target population is the 600 new digital-account customers. The survey misses customers who never open the feature, creating undercoverage. Among feature users, participation is voluntary, so response may be associated with spending or interest in budgeting. Either mechanism can shift the sample mean, and the direction is unknown without more information. More volunteers could reduce random variability within the respondent group while leaving both selection problems in place.

	\secInicio[prob:c9-1][sec:worked-solution-examples]

	\Needspace{5\baselineskip}
	\phantomsection
	\label{sol:c9-1-b}
	\subsubsection*{C9 -- Evaluate the proposed interval and repair the design}

	Even with a simple random sample, $120/600=20\%$ exceeds the usual 10\% guideline for using an independence-based standard error. The actual volunteer sample already lacks random selection. A sample of 120 helps the mean's normal approximation for many populations, but strong skew and an influential extreme value still deserve investigation. Verify the value against a source record and assess its effect on the mean and SD. If it is genuine, keep it in the target data and account for its influence.

	Select distinct customers randomly from the full list, use consistent records or follow up on missing responses, and inspect the spending distribution. A simple random sample of at most 60 meets the sampling-fraction guideline. If more observations are needed, use an analysis that accounts for sampling without replacement, or summarize the full roster if the records are available.

	\secInicio[prob:c9-1][sec:worked-solution-examples]

	\clearpage

	% ----------------------------------------------------------------
	\subsection{Rare defaults in a loan audit (C9)}
	\label{prob:c9-2}

	A lender randomly selects 80 distinct loans from 2,000 loans issued last year. All sampled loans have complete six-month follow-up, and three defaulted in that period. A colleague proposes a normal one-proportion interval because the sample exceeds 30.

	\QuestionsTasks[sec:worked-solution-examples]{
		{9}/{sol:c9-2-a}/
		{Check the sampling fraction and observed outcome counts. Decide whether the usual one-proportion $z$ interval is justified.},
		{9}/{sol:c9-2-b}/
		{State what the lender can report now and how it could obtain a defensible interval.}
	}

	\Needspace{5\baselineskip}
	\phantomsection
	\label{sol:c9-2-a}
	\subsubsection*{C9 -- Check assumptions for this procedure}

	Random selection and $80/2{,}000=4\%$ support approximate independence. The estimate is $\hat p=3/80=0.0375$. There are 3 defaults and 77 non-defaults. The usual one-proportion $z$ interval requires at least 10 of each outcome for its normal approximation, so it is unsuitable here even though $n>30$. The rarity of defaults can leave the sampling distribution of $\hat p$ skewed at this sample size.

	\secInicio[prob:c9-2][sec:worked-solution-examples]

	\Needspace{5\baselineskip}
	\phantomsection
	\label{sol:c9-2-b}
	\subsubsection*{C9 -- Explain a defensible next step}

	Report 3 defaults out of 80 sampled loans, or 3.75\%. If the lender needs an interval from these data, choose a method suited to rare counts and the sampling design. At the observed rate, it would take roughly $10/0.0375\approx267$ loans to expect 10 defaults, which already exceeds 10\% of the 2,000-loan cohort. A larger audit would also need to account for the without-replacement design and check the actual counts.

	\secInicio[prob:c9-2][sec:worked-solution-examples]

	\clearpage

	% ----------------------------------------------------------------
	\subsection{An unusual charge in a small import audit (C9)}
	\label{prob:c9-3}

	An importer randomly draws 12 different freight invoices from 400. Most charges are near 100 USD, but one is 510 USD. All invoices have the same currency and the population SD is unknown.

	\QuestionsTasks[sec:worked-solution-examples]{
		{9}/{sol:c9-3-a}/
		{Does the sampling fraction make a one-sample $t$ interval reliable here? Explain what to do before choosing a method.},
		{9}/{sol:c9-3-b}/
		{A second auditor selects four shipping firms at random and takes three invoices from each firm. Charges from the same firm share a contract rate. Explain what the sampling fraction alone can tell you about independence, and how the auditor should handle this design.}
	}

	\Needspace{5\baselineskip}
	\phantomsection
	\label{sol:c9-3-a}
	\subsubsection*{C9 -- Evaluate the shape requirement for the mean}

	The sampling fraction is $12/400=3\%$, supporting approximate independence. With only 12 values, the extreme charge can strongly affect both $\bar x$ and $s$, and the sample gives reason to question a roughly normal population model. Verify the invoice and the definition of an eligible charge. If the value is genuine and eligible, it belongs in the analysis. A larger random sample or a method appropriate to the skewed cost distribution is preferable to applying the small-sample $t$ interval automatically.

	\secInicio[prob:c9-3][sec:worked-solution-examples]

	\Needspace{5\baselineskip}
	\phantomsection
	\label{sol:c9-3-b}
	\subsubsection*{C9 -- Check dependence within shipping firms}

	Although 12 invoices are only 3\% of the 400, the second design samples in groups. Invoices from one firm can have similar charges because they share a contract rate. The 10\% guideline addresses dependence from sampling without replacement in a simple random sample, while this design also has dependence within firms. Use an analysis that accounts for the grouped design and the selection probabilities, or draw a simple random sample of distinct invoices from the full invoice list.

	\secInicio[prob:c9-3][sec:worked-solution-examples]


\end{document}
